% ==============================================================================
% Fichier : front/03b_labs.tex : FORN (LPI342)
% ==============================================================================

\section*{Les dix labs de NDAP SUKULU}

\begin{tableaucours}{L{2.4cm}L{8.4cm}L{1.1cm}}{%
    \tableauHeader{Lab} &
    \tableauHeader{Sujet} &
    \tableauHeader{Ch.}
}
FORN-LAB-01 & Chaîne de possession et scellés & 5 \\
FORN-LAB-02 & Duplication et vérification par hachage & 8 \\
FORN-LAB-03 & Forensique disque~: lire un système de fichiers & 10 \\
FORN-LAB-04 & Fichiers supprimés et signatures & 11 \\
FORN-LAB-05 & Analyse d'une capture mémoire & 12 \\
FORN-LAB-06 & Forensique réseau~: suivre une exfiltration (flux exportés) & 13 \\
FORN-LAB-07 & Forensique mobile~: lire des artefacts & 14 \\
FORN-LAB-08 & Évaluer la recevabilité d'une preuve & 6 \\
FORN-LAB-09 & Reconstitution chronologique multi-domaines & 17 \\
FORN-LAB-10 & Structurer un rapport forensique & 18 \\
\end{tableaucours}

\medskip
Chaque lab dure environ une heure et se corrige automatiquement. Les pièces
se trouvent dans le dossier \texttt{preuves} (lecture seule) et une copie de
travail dans \texttt{travail}.
